\renewcommand{\myauthor}{Christina Seidner}

\begin{scriptsize}
\begin{lstlisting}[style=Python, language=Python, captionpos=b, stepnumber=5, caption = {backend.py}]
    import os
    import socket
    import logging
    import threading
    import mysql.connector
    from mysql.connector import Error
    from flask import Flask, request, jsonify
    from flask_cors import CORS


    app = Flask(__name__)
    CORS(app)  #activating cors to allow requests from other domains


    #logging-configuration
    logging.basicConfig(level=logging.DEBUG)
    logger = logging.getLogger(__name__)


    #database-configuration
    db_config = {
        "host": os.getenv("DB_HOST", "localhost"),
        "user": os.getenv("DB_USER", "Christina"),
        "password": os.getenv("DB_PASSWORD", "ChrissiNext"),
        "database": os.getenv("DB_NAME", "NextMove"),
    }


    #saving the games
    @app.route('/save_game', methods=['POST'])
    def save_game():
        data = request.json
        user_id = data.get('user_id')
        pgn = data.get('pgn')
        gegner = data.get('gegner')
        ergebnis = data.get('ergebnis')

        try:
            conn = mysql.connector.connect(**db_config)
            cursor = conn.cursor()

            query = """INSERT INTO spiele (user_id, gegner_name, ergebnis, pgn_daten) 
                    VALUES (%s, %s, %s, %s)"""
            cursor.execute(query, (user_id, gegner, ergebnis, pgn))

            conn.commit()
            return jsonify({"status": "Spiel gespeichert!"}), 201
        except Exception as e:
            return jsonify({"error": str(e)}), 500


        finally:
            if conn and conn.is_connected():
                cursor.close()
                conn.close()


    #get games from the database
    @app.route('/get_games/<int:user_id>', methods=['GET'])
    def get_games(user_id):
        conn = None
        try:
            conn = mysql.connector.connect(**db_config)
            cursor = conn.cursor(dictionary=True)

            query = "SELECT * FROM spiele WHERE user_id = %s"
            cursor.execute(query, (user_id,))
            games = cursor.fetchall()
            
            return jsonify({"games": games}), 200
        except Exception as e:
            return jsonify({"error": str(e)}), 500

        finally:
            if conn and conn.is_connected():
                cursor.close()
                conn.close()


    #receive data from board and save games
    def receive_file():
        # Server initialization
        server_address = ("10.5.5.110", 5001)
        server_socket = socket.socket(socket.AF_INET, socket.SOCK_STREAM)
        server_socket.bind(server_address)
        server_socket.listen(1)
        print(f"Server running on {server_address}")

        while True:
            connection, client_address = server_socket.accept()
            print(f"Connection from {client_address} established.")

            file_path = "received_game.pgn"

            try:
                with open(file_path, "wb") as file:
                    while True:
                        data = connection.recv(1024)
                        if not data:
                            break
                        file.write(data)
                print(f"File sucessfully saved to '{file_path}'")
            except Exception as e:
                print(f"Fehler: {e}")
            finally:
                connection.close()


    #main
    if __name__ == "__main__":
        #starting task
        socket_thread =  threading.Thread(target=receive_file, daemon=True)
        socket_thread.start()

        #starting flask app
        app.run(host="10.5.5.110", port=5000)


\end{lstlisting}
\end{scriptsize}